\scanpage{11}
% archive foliation: 10
\uline{Théorie axiomatique des morphismes \unsure{normés|normants}}

Soit $C$ une catégorie, \inserted{\unsure{où les sommes finies existent}}, et soit $\mathcal{N}$ un ensemble de morphismes dans $C$, \deleted{appelés \unsure{morphismes normés}}, ayant les propriétés suivantes

(i) Si $X\in C$, alors $\operatorname{id}_X\in\mathcal{N}$

(ii) Si $u,v\in\mathcal{N}$ sont composables, alors $vu\in\mathcal{N}$

(iii) Si $u:X\longrightarrow Y$ est dans $\mathcal{N}$, et si $Y'\to Y$ est quelconque, alors $X'=X\times_Y Y'$ existe et $u^{Y'}:\deleted{\unsure{X\times_Y}}X'\to Y'$ est $\in\mathcal{N}$

(iv) Si $u_i:X_i\to Y$ $(i=1,2)$ sont $\in\mathcal{N}$, alors $(u_1,u_2):X_1\amalg X_2-Y$ est dans $\mathcal{N}$ ; de plus, l'\unsure{existence} de \unsure{la} \illegible{1} \inserted{$\unsure{\prod}$ et $\unsure{\coprod}$} \unsure{compatible à la formation de} \unsure{sommes}

(v) Si $u:X\to Y$ est $\in\mathcal{N}$, alors \inserted{$\forall n$ \unsure{tout} ss-groupe $G$ de $\unsure{\mathfrak{S}_n|\mathcal{S}_n}$,} le quotient $(\overbrace{X\times_Y\cdots\times_Y X}^{n\text{ fois}})/G$ existe, \inserted{\unsure{et est comp. aux changements de base}}\margin{\unsure{!|?}}

\deleted{\unsure{Il est clair que}}
Soit $\deg$ une fonction sur $\mathcal{N}$, \inserted{notation aussi : $\deg(X/Y)$,} ayant les propriétés suivantes

(i bis) $\deg(\operatorname{id}_X)=1$

(ii bis) $\deg(vu)=\deg v\,\deg u$

(iii bis) $\deg(u)=\deg(u^{Y'})$

(iv bis) $\deg(u,v)=\deg u+\deg v$

(v bis) Le quotient considéré dans (v) a \uuline{\unsure{un}} degré \deleted{\unsure{fini}} \unsure{se calcule} ainsi : Soit \deleted{$\unsure{r=\deg(X/Y)}$}\inserted{$r=\deg(X/Y)$}, $\unsure{A}$ \inserted{$E=[1,r]$},
\deleted{$E=\Delta^r,\quad F=(\mathop{\otimes}\nolimits^r E)/G$ \unsure{deux} \illegible{3}}
Soit $F=E^n/G$, alors \deleted{$\unsure{\deg(X/Y)^n/G}$} le degré de $(\underbrace{X\times_Y-\times_Y X}_{n})/G$ est égal à $\operatorname{card}(E^n/G)$.

Enfin, pour \deleted{\unsure{tout}} morphisme $u\in\mathcal{N}:X-Y$ \inserted{de degré $r$}, \unsure{sur}\ $\mathcal{N}$- soit donné un morphisme
\[
\unsure{\operatorname{Cor}|\operatorname{Cn}}_u=\unsure{\operatorname{Cor}|\operatorname{Cn}}_{X/Y}:Y-\Sigma^r X
\]
(\inserted{appelé} \deleted{\unsure{noté}} \uline{norme} de $u$, ou de $X/Y$), avec les propriétés suivantes

(i ter) $\unsure{\operatorname{Cor}|\operatorname{Cn}}_{X/X}=\operatorname{id}_X$

(ii ter) La commutativité dans le diagramme suivant :
